% Generated draft: requires mathematical and source review.
\documentclass[11pt]{article}
\usepackage[margin=25mm]{geometry}
\usepackage{fontspec}
\setmainfont{Latin Modern Roman}
\usepackage{amsmath,amssymb,mathtools,unicode-math}
\setmathfont{Latin Modern Math}
\usepackage{enumitem,needspace,xurl,hyperref}
\setlength{\parindent}{0pt}
\setlength{\parskip}{5pt}
\setlength{\emergencystretch}{3em}
\newcommand{\N}{\mathbb{N}}
\newcommand{\Z}{\mathbb{Z}}
\newcommand{\Q}{\mathbb{Q}}
\newcommand{\R}{\mathbb{R}}
\newcommand{\C}{\mathbb{C}}
\newcommand{\F}{\mathbb{F}}
\newcommand{\A}{\mathbb{A}}
\newcommand{\PP}{\mathbb{P}}
\newcommand{\E}{\mathbb{E}}
\newcommand{\eps}{\varepsilon}
\newcommand{\dd}{\,\mathrm{d}}
\begin{document}
\section*{UnsolvedMath ID: 2307069}
\textbf{Source problem code:} catalogue entry
\subsection{Definitions and mathematical statement}
Let $A$ be an associative unital algebra with countable basis, and let $V$ be spanned by $e_1,\ldots,e_k$. In coordinates $x_1,\ldots,x_k$, define the left differential operator
\[
Df=\sum_{j=1}^{k}e_j\frac{\partial f}{\partial x_j}.
\]
A differentiable $A$-valued function on $U\subseteq V$ is left analytic in the source's terminology when $Df=0$ throughout $U$. This term does not itself mean ordinary real analyticity. The algebra's scalar field and topology, needed for differentiation and integration of general countably dimensional $A$-valued functions, are not supplied.

The source attributes to Ryan an equivalence between a generalized Cauchy formula with real-analytic kernel and the existence of $p_1,\ldots,p_k\in A$ such that
\[
p_je_l+p_le_j=2\delta_{jl}1\qquad(1\leq j,l\leq k).
\]
The reported formula is
\[
f(x_0)=\int_{\partial M}G(x,x_0)
\left(\sum_{j=1}^{k}(-1)^je_j\,d\widehat x_j\right)f(x),
\]
for left-analytic $f$, an arbitrary compact real $n$-dimensional manifold $M$ lying in $U$, and $x_0\in\mathring M$. Here the conventional coordinate form $d\widehat x_j$ is the wedge product of $dx_1,\ldots,dx_k$ with $dx_j$ omitted. The order of the algebra factors is retained.

\textbf{Questions (J. Ryan).}
\begin{itemize}
\item Does the same algebraic characterization of existence of such a Cauchy formula remain valid if the kernel is required only to be $C^1$, rather than real analytic?
\item What corresponding Cauchy-formula characterization holds for a non-associative algebra?
\end{itemize}
The source does not specify the kernel's domain or permitted singularity at $x=x_0$; its regularity must be distinguished from regularity across that singularity. It also uses dimension $n$ although the displayed boundary form has degree $k-1$, so a literal integral over $\partial M$ needs a compatible dimension and orientation. For the non-associative extension, parenthesization of the triple product and the permitted algebra identities must be fixed. No such missing choices are silently supplied. The status summary reports the question open as of the 2018 edition, with no definitive resolution located.
\subsection{Short English statement}
Does Ryan's algebraic Cauchy-formula criterion survive with a $C^1$ kernel, and what replaces it in non-associative algebras? The analytic framework, kernel singularity, and dimension conventions remain unspecified.
\subsection{Sources}
\begin{itemize}
\item Hayman, \emph{Research Problems in Function Theory} (2018), Problem 7.69, attributed to J. Ryan, \url{https://arxiv.org/abs/1809.07200}; source-TeX extraction, accessed 2026-07-29. The 2018 update was retained for status triage.
\item Ryan, the reported real-analytic-kernel characterization. No separate publication, date, theorem number, or exact analytic hypotheses are supplied.
\item Local registry \texttt{data/unsolvedmath/problems.json}, permanent ID $2307069$, source code catalogue entry, \emph{Research Problems in Function Theory --- Problem 7.69}. The record labels this open and OPEN-TRIAGE, with NEEDS\_REVIEW source-status and rights qualifications.
\end{itemize}
\end{document}
